\newpage

\section*{Обозначения и сокращения}
\addcontentsline{toc}{section}{Обозначения и сокращения}

В работе используется единая система обозначений для рассматриваемых методов.
Часть обозначений относится к транспортной постановке, часть -- к общей выпуклой
и распределенной оптимизации.
Ниже приведено согласование обозначений, которое используется в дальнейшем тексте.

\begin{description}
    \item[$\cX$, $D$] допустимое множество; в транспортной задаче обычно множество путевых потоков, в ML-задаче множество параметров модели.
    \item[$D_1\times\dots\times D_n$] декартова блочная структура допустимого множества.
    \item[$x$] переменная оптимизации в общей задаче и в задачах машинного обучения.
    \item[$f$] поток по ребрам транспортной сети; также $f(x)$ обозначает целевую функцию в общей оптимизации, если это не конфликтует с контекстом.
    \item[$\Psi(f)$] потенциал Бекмана в задаче равновесного распределения транспортных потоков.
    \item[$\tau(f)$] вектор времен движения по ребрам; равен $\nabla\Psi(f)$.
    \item[$\Theta$] матрица принадлежности ребер путям, переводящая путевые потоки в потоки по ребрам.
    \item[$OD$] множество пар источник--сток.
    \item[$\cP_{ij}$] множество путей между источником $i$ и стоком $j$.
    \item[$d_{ij}$] спрос между источником $i$ и стоком $j$.
    \item[$s_k^{FW}$] решение линейного оракула Frank--Wolfe на итерации $k$.
    \item[$d_k$] направление шага; в транспортных главах не смешивается со спросом $d_{ij}$, потому что индекс и контекст различают эти величины.
    \item[$\gamma_k$] длина шага FW.
    \item[$g_k$] зазор Frank--Wolfe.
    \item[$H_k$] гессиан целевой функции в текущей точке.
    \item[$N$] в главах о NFW -- число прошлых направлений; в главе о DBCFW -- число агентов, если явно указан распределенный контекст.
    \item[$B$] размер батча блоков в DBCFW.
    \item[$W_t$] матрица смешивания коммуникационного графа.
    \item[$x_t^i$] локальная точка агента $i$ на итерации $t$.
    \item[$g_t^i$] локальная оценка градиента в градиентном отслеживании.
    \item[$x_{\mathrm{avg},t}$] средняя точка по агентам.
\end{description}

\subsection*{Сокращения}
\addcontentsline{toc}{subsection}{Сокращения}

\begin{description}
    \item[FW] Frank--Wolfe, метод условного градиента.
    \item[CFW] Conjugate Frank--Wolfe.
    \item[BFW] Bi-conjugate Frank--Wolfe.
    \item[NFW] N-conjugate Frank--Wolfe.
    \item[FFW] Fukushima Frank--Wolfe.
    \item[WFFW] Weighted Fukushima Frank--Wolfe.
    \item[BCFW] Block-Coordinate Frank--Wolfe, блочно-координатный Frank--Wolfe.
    \item[SOFW] Stochastic Origin Frank--Wolfe.
    \item[DFW] Decentralized Frank--Wolfe.
    \item[DBCFW] Decentralized Block-Coordinate Frank--Wolfe.
    \item[LMO] linear minimization oracle, линейный оракул.
    \item[BPR] Bureau of Public Roads travel time function.
    \item[OD] origin--destination pair, пара источник--сток.
\end{description}

\subsection*{Линия развития методов}

Рассматриваемая линия методов может быть интерпретирована как последовательное расширение
проекционно-свободной оптимизации по трем исследовательским осям: геометрия направления,
стоимость линейного оракула и распределенная информационная структура.
Первое направление связано с улучшением геометрии шага Frank--Wolfe.
В транспортной задаче это приводит к методам CFW, BFW, NFW, FFW и WFFW, которые используют
память о предыдущих направлениях для уменьшения зигзагообразного поведения траектории.

Второе направление состоит в адаптации сопряженных FW-методов к задачам машинного обучения.
Главная трудность состоит в стоимости вычисления и хранения гессиана, а также в стоимости
точных произведений гессиана на направления.
Поэтому информация о кривизне приближается конечными разностями градиентов.
Так сохраняется идея сопряженности, но алгоритм остается практически применимым.

Третье направление исследует другой вычислительный предел FW -- стоимость полного линейного
оракула.
В крупной транспортной сети полный оракул требует решать задачи кратчайших путей для всех
источников на каждой итерации.
SOFW заменяет полный оракул случайным блочным оракулом и тем самым связывает транспортный
Frank--Wolfe с теорией стохастических блочно-координатных методов.

Четвертое направление обобщает блочную стохастичность на распределенную оптимизацию.
DBCFW соединяет локальные функции агентов, консенсус через матрицу смешивания, отслеживание
градиента и блочный линейный оракул.
Тем самым случайный выбор блоков становится частью децентрализованной схемы условного
градиента.
